% File example.tex
% Contact: simonnet@ecole.ensicaen.fr
%
% version 1.0 (July 24, 2009)
% version 1.1 (September 3, 2009)
% add using of the optional command: \secondauthoraddr

\documentclass[10pt]{article}

\input{../template/preamble.tex}

%other package

\usepackage{lmodern}
\usepackage{graphicx}
\usepackage{times}
\usepackage{amsmath}
\usepackage{amssymb}
\usepackage{booktabs}
\usepackage[colorlinks=true, urlcolor=blue, linkcolor=black]{hyperref}

\begin{document}
\noindent

% This should produce references in the order they appear
\bibliographystyle{ieeetr}

\title{TP ICD - Análise Estatística da Associação entre Financiamento Educacional, Ineficiência Escolar e Criminalidade no Brasil}

\authorname{Gabriel R. Moura, Mateus M. T. Costa, Pedro R. F. de Aguiar, Renato V. M. P. Pinto}
\authoraddr{Universidade Federal de Minas Gerais (UFMG) - Departamento de Ciência da Computação (DCC)}

\maketitle

\keywords
Criminalidade, Financiamento Educacional, Testes de Permutação, Regressão Múltipla, Bootstrap, K-Means.

\abstract
Este estudo investiga a relação empírica entre os gastos estatais em educação, o desempenho de retenção escolar e as taxas de violência letal (CVLI) nos estados brasileiros. Como a natureza dos dados é observacional, a pesquisa não busca estabelecer relações causais estritas, mas sim identificar e quantificar associações estruturais robustas por meio de inferência estatística. Utilizando testes de hipótese, métodos de reamostragem (Bootstrap) para construção de Intervalos de Confiança, modelos preditivos de Regressão Múltipla e o modelo de classificação de K-Means, fizemos análises de comparação e agrupamento de Estados Brasileiros e estudos da associação de fatores educacionais com a criminalidade. Os resultados apontam que, estatisticamente, maiores aportes financeiros per capita e menores taxas de abandono escolar estão significativamente associados a menores índices de criminalidade letal.

\section{Introdução}

A formulação de políticas públicas baseadas em evidências exige uma transição de heurísticas sociológicas para validações estatísticas rigorosas. O presente projeto propõe uma exploração da interseção entre o aparato educacional estatal (mensurado sob as óticas de aporte financeiro e eficiência de fluxo escolar) e a incidência de criminalidade letal nas Unidades Federativas do Brasil (UF). 

É imprescindível ressaltar que a complexidade do fenômeno da violência impede atribuições causais diretas a partir de dados puramente observacionais. Desse modo, nosso objetivo analítico é dimensionar o poder explicativo dessas covariáveis educacionais sobre a variância da criminalidade. Para nortear a investigação, formalizamos as seguintes Perguntas de Pesquisa:
\begin{enumerate}
    \item \textbf{Pergunta 1:} Como evoluíram as taxas de evasão escolar e de criminalidade ao longo dos anos nos diferentes estados brasileiros?
    \item \textbf{Pergunta 2:} Existe diferença significativa na criminalidade letal entre os estados que mais investem em educação e os que menos investem?
    \item \textbf{Pergunta 3:} Existe correlação entre o investimento financeiro estatal, a ineficiência escolar (reprovação/abandono) e a criminalidade?
    \item \textbf{Pergunta 4:} É possível agrupar os Estados Brasileiros em perfis socioeconômicos para analisar semelhanças e diferenças nas tendências criminais e educacionais das UFs?
\end{enumerate}

\section{Metodologia}

\subsection{Bases de Dados}
A análise baseia-se na integração de dados públicos extraídos de instituições governamentais referentes a instâncias estaduais. As fontes incluem:
\begin{itemize}
    \item \textbf{\href{https://siconfi.tesouro.gov.br/siconfi/pages/public/consulta_finbra_rreo/finbra_rreo_list.jsf}{SICONFI - Tesouro Nacional}:} Despesas consolidadas estaduais na função Educação.
    \item \textbf{Instituto Nacional de Estudos e Pesquisas Educacionais Anísio Teixeira - } \textbf{\href{https://www.gov.br/inep/pt-br/acesso-a-informacao/dados-abertos/indicadores-educacionais/taxas-de-rendimento-escolar}{INEP}:} Indicadores de fluxo escolar, com ênfase nas taxas de aprovação, reprovação e evasão (abandono) no Ensino Médio.
    \item \textbf{\href{https://dados.mj.gov.br/dataset/sistema-nacional-de-estatisticas-de-seguranca-publica}{SINESP - Ministério da Justiça}:} Ocorrências criminais. Isolamos a métrica CVLI (Crimes Violentos Letais Intencionais: homicídios, latrocínios e lesões seguidas de morte) para evitar vieses de subnotificação e de crimes patrimoniais não violentos ligados a estados com maior densidade de riqueza urbana.
    \item \textbf{Instituto Brasileiro de Geografia e Estatistica - } \textbf{\href{https://www.ibge.gov.br/}{IBGE}:} Estimativas populacionais para normalização.
\end{itemize}

\subsection{Fluxo de Trabalho e Tratamento Metodológico}
Para garantir a equivalência paramétrica entre estados heterogêneos, converteu-se o financiamento bruto em Investimento por Aluno e as ocorrências criminais em taxas proporcionais (CVLI por 100 mil habitantes). Além disso, devido à disparidade da faixa de tempo disponível em cada Base, a leitura conjunta ficou restrita ao intervalo \textbf{2015-2022}.

A \textbf{Pergunta 1} correspondeu ao processo de \textbf{Análise Experimental} que se deu por meio da montagem de gráficos e tratamento dos dados para os indíces de evasão escolar e de criminalidade de modo a perceber como os dados se comportavam e obter insigths iniciais.
A \textbf{Pergunta 2} foi abordada via \textbf{Testes A/B}, dividindo os estados de acordo com a mediana de Investimento em Educação, construindo a distribuição nula com permutação, Intervalo de Confiança com Bootstrap e avaliando a significância da diferença das médias de criminalidade entre os grupos, com validação do poder estatístico ($1 - \beta$). A \textbf{Pergunta 3} consistiu na construção de uma \textbf{Regressão Linear Múltipla} otimizada por Gradiente Descendente. O treinamento incluiu transformações logarítmicas, diagnósticos de multicolinearidade (construção de matrizes de correlação) e \textbf{Validação Cruzada K-Fold}, garantindo a qualidade de generalização. A inferência paramétrica dos pesos $\theta$ e da qualidade do modelo mensurada pela métrica $R^{2}$ deu-se de forma empírica via \textbf{Bootstrap}. A \textbf{Pergunta 4}... [DESCREVA AQUI MINHA GENTE POGGERS]

\subsection{Divisão de Tarefas}
Além das atividades comuns para todos os membros do grupo, como a pesquisa das Bases de Dados e tratamento inicial, dividimos a execução da resolução das perguntas da seguinte forma:
\begin{itemize}
    \item \textbf{Gabriel:} Implementação e análise final da EDA, design inicial do processo de clusterização.
    \item \textbf{Mateus:}  Implementação e análise da Pergunta 2 e construção dos Intervalos de Confiança paramêtricos para a Pergunta 3, além da validação cruzada analisada para a Regressão Múltipla.
    \item \textbf{Pedro:}  Elaboração de determinados gráficos para Pergunta 1 e implementação do Pré-Processamento e da Regressão Múltipla da Pergunta 3.
    \item \textbf{Renato:}
\end{itemize}

\section{Desenvolvimento}

\subsection{Caracterização e Análise Exploratória dos Dados}
Para conseguirmos uma compreensão inicial dos dados, investigamos a evolução temporal das taxas de evasão escolar e de criminalidade letal nos estados brasileiros. Essa etapa foi feita no notebook da Pergunta 1 a partir de três bases principais: a base de criminalidade, a base de indicadores escolares e a base com informações populacionais e regionais das UFs. Antes das visualizações, realizamos o pré-processamento necessário para tornar as variáveis comparáveis. Na base escolar, colunas numéricas que estavam armazenadas como texto, como \textit{NumEvasões}, \textit{num\_matricula\_EM} e \textit{tx\_abandono\_EM}, foram padronizadas e convertidas para formato numérico. A variável educacional usada na análise foi a \textit{Taxa\_Abandono\_EM}, derivada diretamente da taxa de abandono do Ensino Médio.

Como as Unidades Federativas possuem populações muito distintas, decidimos normalizar as ocorrências criminais absolutas em taxas por 100 mil habitantes. Para isso, integramos a base de violência à base de estados, adicionando população e região a cada UF, e calculamos $\text{Taxa}_{100k} = \frac{\text{Ocorrências}}{\text{População}} \times 100.000$ para cada tipo de crime. Essa normalização evita que estados muito populosos pareçam mais violentos apenas por terem mais habitantes. 

Além disso, no quesito criminal, usamos como indicador a quantidade de \textbf{CVLIs} (Crimes Violentos Letais Intencionais), construída pela soma de homicídios dolosos, latrocínios e lesões corporais seguidas de morte, uma vez que isso reduz o peso de crimes patrimoniais e que, consequentemente, podem estar mais correlacionados a outros fatores, como urbanização, frota de veículos e nível de atividade econômica. 

Por fim, para a dimensão educacional, utilizamos a taxa de abandono no Ensino Médio, já que essa representa uma etapa crítica de permanência escolar e de transição para o mercado de trabalho. A comparação conjunta ficou restrita ao intervalo \textbf{2015-2022}, período comum entre as bases analisadas.

Para facilitar uma análise visual, construímos gráficos de linha por UF para acompanhar a evolução temporal de cada crime e da taxa de abandono escolar, o que permitiu observar trajetórias individuais, oscilações anuais e possíveis mudanças bruscas em estados específicos. Em seguida, usamos \textbf{mapas de calor (heatmaps)} no formato UF $\times$ ano, uma vez que, embora os gráficos de linha sejam bons para analisar cada estado separadamente, dificultam uma comparação simultânea de todas as UFs. Já nos heatmaps, cada linha representa um estado, cada coluna representa um ano e a intensidade da cor representa a taxa observada, o que permite rapidamente identificar estados com níveis persistentemente altos, anos de pico ou queda, outliers e padrões regionais. Os estados foram ordenados pela média histórica do indicador, de modo que aqueles com maiores valores médios aparecessem no topo.

Foram gerados heatmaps para os indicadores criminais e para a evasão escolar. No caso da criminalidade, eles ajudaram a visualizar a concentração territorial da violência e a distinguir estados com taxas historicamente elevadas daqueles com níveis menores ou em queda. No caso da evasão, o heatmap permitiu identificar estados com abandono escolar estruturalmente alto e estados que apresentaram melhora ao longo da série. Por fim, para relacionar diretamente as duas dimensões da pergunta, calculamos a média estadual de abandono no Ensino Médio e a média estadual de CVLI no período 2015-2022 e construímos um gráfico de dispersão, com cada ponto representando uma UF e a cor indicando sua região geográfica.

A primeira leitura exploratória indicou queda agregada nos dois indicadores. A média estadual de CVLI passou de \textbf{31,37 crimes letais por 100 mil habitantes em 2015} para \textbf{24,24 em 2022}, redução de aproximadamente \textbf{22,7\%}. No mesmo período, a taxa média de abandono no Ensino Médio caiu de \textbf{9,72\%} para \textbf{6,84\%}, o que corresponde a uma redução aproximada de \textbf{29,7\%}. Portanto, no agregado nacional por UF, tanto a violência letal quanto a evasão escolar apresentaram melhora ao longo do tempo.

Apesar dessa tendência geral de queda, o cenário não é homogêneo entre os estados. Na criminalidade, as maiores médias de CVLI entre 2015 e 2022 ocorreram no \textbf{Rio Grande do Norte (46,92)}, \textbf{Alagoas (43,80)}, \textbf{Sergipe (42,67)}, \textbf{Pernambuco (42,01)} e \textbf{Ceará (40,02)}. Os menores valores médios foram observados em \textbf{São Paulo (7,70)}, \textbf{Santa Catarina (10,40)}, \textbf{Minas Gerais (15,64)}, \textbf{Paraná (15,70)} e \textbf{Distrito Federal (16,05)}. Esses dados demonstram um possível componente regional na concentração de casos de violência letal, estando esses localizados especialmente em estados do Norte e do Nordeste.

Na evasão escolar, também houve forte variação entre UFs. As maiores médias de abandono no Ensino Médio foram registradas no \textbf{Pará (12,62\%)}, \textbf{Rio Grande do Norte (9,93\%)}, \textbf{Paraíba (9,49\%)}, \textbf{Mato Grosso (9,30\%)} e \textbf{Alagoas (9,26\%)}. Por outro lado, os menores valores médios apareceram em \textbf{Pernambuco (1,49\%)}, \textbf{Espírito Santo (3,15\%)}, \textbf{Goiás (3,23\%)}, \textbf{São Paulo (3,30\%)} e \textbf{Distrito Federal (4,01\%)}. Assim como na criminalidade, a melhora agregada convive com desigualdades territoriais, e a questão regional se apresenta novamente.

Por fim, comparamos as médias estaduais de abandono no Ensino Médio e de CVLI no período 2015-2022. A correlação encontrada foi de \textbf{$r = 0,34$}, indicando uma associação positiva moderada: estados com maior abandono escolar tendem, em média, a apresentar maior violência letal. No entanto, a dispersão entre os estados mostra que a criminalidade não pode ser explicada apenas pela evasão escolar. Fatores como desigualdade de renda, urbanização, capacidade institucional, políticas de segurança, mercado de trabalho e dinâmicas locais de violência também influenciam o fenômeno.

Dessa forma, a resposta exploratória à Pergunta 1 é que as taxas de evasão escolar e de criminalidade letal caíram no agregado entre 2015 e 2022, mas de maneira desigual entre os estados. A relação entre abandono escolar e CVLI aparece como uma evidência estatística de associação, suficiente para orientar as análises posteriores, mas não como demonstração direta de causalidade.

\subsection{Teste de Hipóteses (Pergunta 2: Efeito do Financiamento na Violência)}

Para consulta da implementação e da visualização das distribuições construídas, acesse o arquivo \href{https://colab.research.google.com/drive/1dT2fBm92cKXEZFFhQeh7Xt-DvyE9x9aO#scrollTo=7e053f8b-1a88-4c0c-a0c7-7ba926015c9e}{Pergunta2.ipynb}.

Para avaliar de forma probabilística se a disparidade de recursos per capita implica em reduções sistêmicas da letalidade criminal, segregamos a base de dados em dois subgrupos: Estados de Alto Investimento (acima da mediana): \textbf{Grupo A}; e Estados de Baixo Investimento (abaixo da mediana): \textbf{Grupo B}.

A literatura sociológica e o senso comum sugerem que a precarização da educação piora a segurança. Logo, nossa pergunta não era apenas se os grupos eram "diferentes", mas especificamente se o Grupo B (baixo investimento) possuía uma taxa de criminalidade maior que o Grupo A. Formalizamos as hipóteses matematicamente como:
\begin{itemize}
    \item \textbf{Hipótese Nula ($H_{0}$):} $\mu_{B} - \mu_{A} = 0$
    \item \textbf{Hipótese Alternativa ($H_{A}$):} $\mu_{B} - \mu_{A} > 0$
\end{itemize}

Sendo $\mu_{A} = \text{Média de CLVI do Grupo A}$ e $\mu_{B} = \text{Média de CLVI do Grupo B}$.

Como métricas como CVLI e financiamento estatal raramente seguem distribuições normais perfeitas (\textit{outliers} e assimetria alta), além de maior coerência com o foco em abordagens computacionais da disciplina, adotou-se o \textbf{Teste de Permutação}, um método não-paramétrico de reamostragem sob a Hipótese Nula ($H_0$) de que os rótulos "Alto" e "Baixo" investimento não alteram a distribuição das taxas de criminalidade. Além disso, para construção do Intervalo de Confiança para a média de criminalidade, utilizamos outro método não-paramêtrico do \textbf{Bootstrap}.

\paragraph{Resultados e Poder do Teste:}
Os dados revelaram uma diferença média observada de 8.91 CVLIs por 100.000 habitantes em desfavor dos estados com menor investimento educacional. A metodologia empírica gerou um P-valor aproximado de $0.013$ ($1.33\%$). Sob o nível de significância crítico $\alpha = 0.05$, rejeitamos $H_0$ ($0.013 < 0.05$). Ou seja, com um nível de confiança de (95\%), os dados amostrais fornecem evidências suficientes para rejeitar a hipótese nula ($H_{0}$) em favor da hipótese alternativa ($H_{A}$), levando à conclução de que os estados brasileiros que aplicam menos recursos proporcionais na educação possuem taxas sistematicamente maiores de crimes letais. 

Mais fundamental do que apontar a significância é avaliar a consistência deste teste. Computamos empiricamente o erro do Tipo II ($\beta$) observando a sobreposição do cenário nulo com a distribuição amostral simulada via Bootstrap. Obtivemos $\beta = 0.2563$. Logo, a \textbf{Força do Teste ($1 - \beta$) é estimada em $74.37\%$}. Considerando a magnitude da disparidade das ocorrências criminais letais associada ao baixo investimento em educação (diferença média observada de 8.91 e P-Valor de 1.33\%), se essa associação estrutural for a realidade do país, o nosso modelo possui 74.37\% de probabilidade de identificar essa deficiência socioeducacional com sucesso. Esse nível de poder estatístico (próximo à referência acadêmica comum de 80\%) indica que o teste tem uma sensibilidade robusta para detectar o efeito.

\subsection{Modelagem de Regressão Múltipla (Pergunta 3: Financiamento, Ineficiência Escolar x Criminalidade) }

Para consulta da implementação e da visualização do Modelo e dos resultados obtidos, acesse o arquivo \href{https://colab.research.google.com/drive/1pVdoor4xno_OC3BWHZh3yFA97_0DO5bY}{Pergunta3.ipynb}.

Desenvolveu-se um modelo de Regressão Linear Múltipla, focado em mensurar simultaneamente os efeitos marginais isolados do \textit{Investimento por Aluno} ($X_1$) e das variáveis de ineficiência escolar \textit{Taxa de Abandono} ($X_2$) e \textit{Taxa de Reprovação} ($X_3$) sobre a criminalidade letal ($Y$). A função de custo (Erro Quadrático Médio) foi minimizada através do algoritmo de \textbf{Gradiente Descendente}.

\paragraph{Engenharia de Atributos e Diagnóstico de Multicolinearidade:}
A fase inicial da modelagem revelou a necessidade de ajustes algébricos e estatísticos vitais:
\begin{enumerate}
    \item \textit{Visualização de Não-Linearidade:} O gráfico de dispersão do Investimento revelou rendimentos marginais decrescentes. Aplicou-se a transformação $Log(X_1)$, melhorando o ajuste da reta ao espaço de atributos originais.
    \item \textit{Rigor contra Multicolinearidade:} Uma observação superficial apontaria para a inclusão conjunta das variáveis de evasão e reprovação. Contudo, a matriz de correlação cruzada revelou que a Taxa de Reprovação possuía correlação desprezível com a variável alvo, CVLI ($\rho \approx 0.09$), e alta colinearidade (Pearson = $0.47$) com o Abandono. Na teoria de Regressão Múltipla, introduzir variáveis correlacionadas entre si (Multicolinearidade) e que não agregam poder explicativo ao alvo pode gerar instabilidade nos pesos do modelo e aumentar a variância sem ganho real de predição. Pelo Princípio da Parcimônia, optamos por remover a Taxa de Reprovação das modelagens seguintes, construindo um modelo mais simples, estável e interpretável focado apenas em Investimento e Abandono.
\end{enumerate}

\paragraph{Validação da Generalização (Prevenção de \textit{Data Leakage}):}
Para avaliar o poder explicativo do modelo reduzido ($Y \sim \theta_0 + \theta_1 Log(Investimento) + \theta_2 Abandono$), aplicou-se a \textbf{Validação Cruzada Estratificada (5-Fold CV)}. 

Em coerência aos preceitos estritos da aprendizagem de máquina estatística, a normalização dos dados (\textit{Z-Score}) foi processada de forma isolada nos dados de treino, isto é, média e desvio padrão foram computados exclusivamente nos dados do \textit{fold} de treinamento e posteriormente projetados sobre o teste. Esse rigor impediu o vazamento de informações (\textit{Data Leakage}). Obtivemos um $R^2$ médio de teste de $0.2175$ ($\pm 0.0615$), muito próximo ao de treino ($0.2407$), refutando a presença de superajuste (\textit{overfitting}) e atestando a capacidade de generalização da arquitetura estabelecida.

\paragraph{Inferência por Bootstrap e Intervalos de Confiança:}
Para inferência final, evitando o viés otimista que ocorreria caso selecionássemos arbitrariamente o "melhor" modelo do \textit{K-Fold}, o vetor paramétrico final foi treinado aproveitando a amostra agregada total. Empregou-se o \textbf{Bootstrap} (4.000 reamostragens com reposição) para estimar a densidade probabilística dos coeficientes.

Os Intervalos de Confiança (IC a 95\%) revelaram as seguintes associações puras:
\begin{itemize}
    \item \textbf{Coef. Log(Investimento) IC:[-6.0029, -3.2036]:} Isso indica com robustez estatística que, independentemente da ineficiência escolar, maiores investimentos estatais por aluno estão diretamente associados a menores taxas de violência (CVLI). O sucesso da transformação logarítmica sugere ainda que a relação não é perfeitamente linear: os maiores impactos sociais e reduções de criminalidade ocorrem nos incrementos iniciais de investimento nos estados historicamente mais carentes.
    \item \textbf{Coef. Abandono IC:[1.5027, 4.4105]:} Isso comprova matematicamente que a evasão escolar (o momento em que o jovem se desvincula do sistema educacional e fica desassistido pelo Estado) está de forma ativa e isolada relacionada com a escalada da violência.
    \item \textbf{Qualidade de Ajuste ($R^2$) IC:[15,14\% a 33,86\%]:} Embora existam inúmeros outros fatores socioeconômicos que afetem a violência, provamos estatisticamente que apenas duas dimensões (financiamento e evasão escolar) são capazes de explicar, de forma consistente, cerca de um quinto a um terço de toda a variabilidade da criminalidade entre os estados, com 95\% de confiança.
\end{itemize}

\paragraph{Veredito:}
Sim, a correlação existe e foi validada com rigor estatístico. Contudo, os dados revelam que "ineficiência escolar" não é um bloco homogêneo. As políticas públicas com maior potencial de retorno na segurança pública estadual devem focar em duas frentes simultâneas e independentes: a injeção bruta de recursos financeiros (Investimento per capita) e a criação de mecanismos de retenção para evitar a evasão (Abandono). O combate exclusivo à taxa de reprovação, isoladamente, mostrou-se estatisticamente irrelevante para a mitigação da criminalidade.

\subsection{Modelo de Classificação/Clustering (Pergunta 4)}

MAIS PAPINHO BLA BLAH

\section{Conclusões}

A investigação evidenciou de forma quantificável que o aparelho educacional atua fortemente na macroestrutura da segurança pública brasileira. Com resguardo da inferência não-causal decorrente de dados estritamente observacionais, foi possível validar estruturalmente a correlação proposta nas hipóteses iniciais.

No âmbito financeiro, tanto os testes de hipóteses de permutação quanto a análise marginal da Regressão Múltipla convergem para o mesmo fato estatístico: a letalidade criminal apresenta redução considerável em cenários de aportes governamentais crescentes.

Ademais, a decomposição empírica da ineficiência escolar forneceu um achado sociológico vital: o impacto pró-criminalidade reside especificamente na evasão escolar (Abandono). Os dados revelaram que a mera retenção do aluno na mesma série (Reprovação), desde que ele permaneça dentro da infraestrutura da escola, não é um motor significativo para a criminalidade.

\end{document}
